\documentclass[../../main.tex]{subfiles}
\begin{document}

\begin{center}
  {\fontsize{48pt}{58pt}\selectfont \textbf{1988}\normalsize\textbf{年度}}
\end{center}

\begin{table}[htb]
  \centering
  \caption{出題テーマと難易度}
  \begin{tabular}{|c|p{9cm}|c|}
    \hline
    \textbf{問題番号} & \textbf{テーマ} & \textbf{難易度} \\
    \hline
    第1問 & 数列の極限 & \\
    \hline
    第2問 & 二次関数と導関数の評価 & \\
    \hline
    第3問 & 三角関数のグラフと円の交点 & \\
    \hline
    第4問 & 円周上の動点間の距離 & \\
    \hline
    第5問 & 数列の極限（二項係数） & \\
    \hline
  \end{tabular}
\end{table}

\begin{table}[htb]
  \centering
  \caption{解答}
  \renewcommand{\arraystretch}{1.4}
  \begin{tabular}{|c|c|p{9cm}|}
    \hline
    \textbf{問題番号} & \textbf{小問} & \textbf{答え} \\
    \hline
    第1問 & - & $\displaystyle\lim_{n\to\infty}a_n=1$ \\
    \hline
    第2問 & (1) & 証明問題（$|f'(1)|\le4$） \\
    \cline{2-3}
    & (2) & $(a,b,c)=(\pm2,0,\mp1)$，すなわち $f(x)=\pm(2x^2-1)$ \\
    \hline
    第3問 & - & $N(r)$: $0<r<\sqrt{\dfrac{\pi}{8}+\dfrac{1}{2}}$で$0$個，$r=\sqrt{\dfrac{\pi}{8}+\dfrac{1}{2}}$または$1<r$で$2$個，$r=1$で$3$個，$\sqrt{\dfrac{\pi}{8}+\dfrac{1}{2}}<r<1$で$4$個 \\
    \hline
    第4問 & - & 最大値 $\dfrac{5\sqrt{5}}{2}$，最小値 $1$ \\
    \hline
    第5問 & - & $\dfrac{27}{16}$ \\
    \hline
  \end{tabular}
\end{table}

\paragraph{試験形態}
5問．試験時間は要確認．

\end{document}
